\clearpage\thispagestyle{otFront}\phantomsection\label{nouveautes}\pdfbookmark[0]{Ce qui change}{bm-nouveautes}\otPageTitre{Ce qui change}{Nouveau programme de Quatrième · Rentrée 2027}\fontsize{11}{15}\selectfont

Le programme de mathématiques du cycle 4 publié au Bulletin officiel du 5 mars
2026 entre en vigueur en Quatrième à la rentrée 2027. Cette édition a été
entièrement reprise pour le suivre.\par\medskip
\begin{savoirfaire}[Ce qui entre en Quatrième]Nombres rationnels : définition, inverse, fraction d’une fraction. Puissances : $a^m\times a^n=a^{m+n}$ et $a^n\times b^n=(ab)^n$. Ratios et partages proportionnels ; coefficient multiplicateur ; programmes de calcul que l’on remonte et notation $P(t)$. Étendue, effet d’une valeur extrême et tableur. Notations ensemblistes $\overline A$, $A\cap B$, $A\cup B$, $\varnothing$ ; expériences à deux épreuves ; fluctuation des fréquences. Droites remarquables, réciproque, contraposée et raisonnement par l’absurde ; droite des milieux ; triangle rectangle et cercle circonscrit ; rectangles sans équerre. Conditions avec $<$, $>$, $=$ et une variable.\end{savoirfaire}
\begin{attention}[Ce qui passe en Troisième]Théorème de Thalès ; cosinus, sinus et tangente ; puissances d’exposant négatif et écriture scientifique ; décomposition en facteurs premiers et fractions irréductibles ; conditions composées, boucles « répéter jusqu’à » et programmes à plusieurs variables.\end{attention}
\begin{exemple}[Ce qui n’est plus au programme de Quatrième]Cas d’égalité des triangles ; agrandissements et réductions avec leurs effets sur les aires et les volumes ; repérage dans un pavé droit ; grandeurs produits ; probabilités calculées par des rapports d’aires. Les sommes de fractions de dénominateurs quelconques, les diagrammes circulaires et les volumes du prisme et du cylindre relèvent de la Cinquième dans le programme 2026 : ils sont réactivés ici comme automatismes.\end{exemple}
Le repérage détaillé, objectif par objectif, figure à la fin du livre.

\clearpage
